% ###################
% ## Filename: sCM_inventory.tex
% ## Location: D:\Project\Prej_Teaching\ISE\SupplyChainMgt\presentation\scm_05_aggreated
% ## Created: Jan 26, 2011
% ## Last updated: Feb 2, 2010
% ## Figure Location: ..\Figures\
% ## Note: Presentation, Teaching
% ###################

%\documentclass[10pt,presentation]{beamer}
%\documentclass[10pt,handout]{beamer}
\documentclass[10pt,handout,notes=show]{beamer}

%============================================================================================================--
\usepackage{fontspec}
\usepackage{xunicode}
\usepackage{xltxtra}
\usepackage{array}
\usepackage{xcolor}
\usepackage{verbatim}
\usepackage{beamerthemesplit}
\usepackage{beamerbasetemplates}
\usepackage{graphics}
\usepackage{graphicx}
\usepackage{hyperref}
\usepackage{multirow}
\usepackage{tikz,pgfplots}
\usepackage{algorithm}
\usepackage{algorithmic}
\usepackage{soul,xcolor}
\usepackage{colortbl}
\usepackage{multicol}

\XeTeXlinebreaklocale"th"
\defaultfontfeatures{Scale=1.5, Mapping=tex-text} %{Scale=1.23}



\usetheme{Madrid}  % JuanLesPins  Rochester
           %  Berkeley  Palo Alto    with sidebar and top
           % Goettingen  Marburg     with sidebar
           %  Copenhagen  Luebeck  Warsaw

\setbeamertemplate{headline}
{%
    \begin{beamercolorbox}{section in head/foot}
        \vskip-2pt \hskip5pt
        \insertsectionnavigationhorizontal{\textwidth}{}{}
        \hskip5pt \vskip-2pt
    \end{beamercolorbox}
}%
\setbeamertemplate{footline}
{%
    \begin{beamercolorbox}{section in head/foot}
        \vskip-2pt \hskip5pt
        \insertshorttitle \hfill \insertshortdate \hskip10pt \insertframenumber / \inserttotalframenumber
        \vskip-2pt\hskip5pt
    \end{beamercolorbox}
}%
\definecolor{alertedOrange}{RGB}{255,36,0}
\setbeamercolor{alerted text}{fg=alertedOrange}
\useinnertheme{rounded}
\usecolortheme{rose}
%\setbeamercovered{transparent}
\usefonttheme{structuresmallcapsserif}
\definecolor{LRed}{rgb} {1,.8,.8}
\definecolor{LBlue}{rgb}{.8,.8,1}

%%%%%%%%%%%%%%%%%%%%%%%%%%%%%%%%%%%%%%%%%%%%%%%%%%%%%%%%%%%%%%%%%%%%%%%%%%%%%%%
% Notation:
\newtheorem{observation}{Observation}
\newcommand{\figureLocation}[1]{../figures/#1}
\newcommand{\figureLoc}[1]{../figures/#1}
\newcommand{\refLocation}[1]{../../../../Reference/#1}
\newcommand{\NextLine}{\vskip 0.05\textheight}
\newcommand{\NextHalfLine}{\vskip 0.02\textheight}
\newcommand{\RemoveLine}{\vskip -0.1\textheight}
\newcommand{\RemoveHalfLine}{\vskip -0.02\textheight}
\newcommand{\BackTab}{\hskip -0.15\textwidth}
\newcommand{\HalfBackTab}{\hskip -0.11\textwidth}
\newcommand{\DualFigureSize}{1.10\textwidth}
\newcommand{\MiniPageSize}{0.52\textwidth}
\newcommand{\code}[1]{\textcolor{blue}{#1} }
\newcommand{\rightNote}[2]{\hfill \uncover<#1>{\textbf{\code{\#\# {#2}  }}}}
\newcolumntype{P}[1]{>{\raggedright\arraybackslash}p{#1}<{}}

\author[Kittithreerapronchai]{%
  Oran~Kittithreerapronchai\inst{1} }
\institute[IE Chula]{
  \inst{1}%
  Department of Industrial Engineering, Chulalongkorn University\\
  Bangkok 10330 THAILAND \\
  }
\subject{}
\date{}

\setbeamertemplate{enumerate items}[default]
\setbeamertemplate{bibliography item}[text]
\bibliographystyle{alpha}
\newcommand{\textFont}[1]{\footnotesize{#1}}
\setcounter{MaxMatrixCols}{20}

%%%--------------------------------

\makeatletter
\defbeamertemplate{note page}{infolines}
{%
  {%
    \small
    \insertvrule{.25\paperheight}{white!90!black}
    \vskip-.25\paperheight
    \nointerlineskip
    \vbox{
      \hfill\insertslideintonotes{0.25}\hskip-\Gm@rmargin\hskip0pt%
      \vskip-0.25\paperheight%
      \nointerlineskip
      \begin{pgfpicture}{0cm}{0cm}{0cm}{0cm}
        \begin{pgflowlevelscope}{\pgftransformrotate{90}}
          {\pgftransformshift{\pgfpoint{-2cm}{0.2cm}}%
          \pgftext[base,left]{\footnotesize\the\year-\ifnum\month<10\relax0\fi\the\month-\ifnum\day<10\relax0\fi\the\day}}
        \end{pgflowlevelscope}
      \end{pgfpicture}}
    \nointerlineskip
    \vbox to .25\paperheight{\vskip0.5em
      \hbox{\insertshorttitle[width=8cm]}%
      \setbox\beamer@tempbox=\hbox{\insertsection}%
      \hbox{\ifdim\wd\beamer@tempbox>1pt{\hskip4pt\raise3pt\hbox{\vrule
            width0.4pt height7pt\vrule width 9pt
            height0.4pt}}\hskip1pt\hbox{\begin{minipage}[t]{7.5cm}\def\breakhere{}\insertsection\end{minipage}}\fi%
      }%
      \setbox\beamer@tempbox=\hbox{\insertsubsection}%
      \hbox{\ifdim\wd\beamer@tempbox>1pt{\hskip17.4pt\raise3pt\hbox{\vrule
            width0.4pt height7pt\vrule width 9pt
            height0.4pt}}\hskip1pt\hbox{\begin{minipage}[t]{7.5cm}\def\breakhere{}\insertsubsection\end{minipage}}\fi%
      }%
      \setbox\beamer@tempbox=\hbox{\insertshortframetitle}%
      \hbox{\ifdim\wd\beamer@tempbox>1pt{\hskip30.8pt\raise3pt\hbox{\vrule
            width0.4pt height7pt\vrule width 9pt
            height0.4pt}}\hskip1pt\hbox{\insertshortframetitle[width=7cm]}\fi%
      }%
      \vfil}%
  }%
  \vskip.25em
  \nointerlineskip
  %\setmainfont{TH Sarabun New}
  \begin{minipage}{\textwidth} % this is an addition
  \footnotesize{\fontspec{TH Sarabun New} \insertnote}
  \end{minipage}               % this is an addition
  %\setmainfont{Computer Modern}

}
\makeatother

\setbeamertemplate{note page}[infolines]

\title[LSCM v3.0]{Lecture 04:\\ Demand Management and \\ Aggregate Planning }
\hypersetup{
  pdftitle={04 Demand Management and Aggregate Planning},
  pdfsubject={2104509lscm & 2184408scm handout},
  pdfproducer={XeLatex/WinEdt},
  pdfcreator={oran.k@chula.ac.th}
}


\begin{document}

%--------------------------------------------------------------------------------------------------------------
\begin{frame}[t]
  \titlepage \label{Slide:Title}

  \NextLine   \NextLine
  \begin{flushright}
    \tiny{ \textbf{last updated:} \today } \hskip3ex
  \end{flushright}

  \note{
    \textbf{Purpose:} introducing SCM practical question and focusing on demand planning and aggregate planning\\
    \textbf{Leading Question:} what are SCM questions that a company may encounter and what are their related disciplines\\
    \begin{itemize}
      \item many SCM textbook overlapped with Operation Management textbook. Why?
      \item How operations management, a study of operational activities in a company, related to SCM?
      \item What are the most valuable information in SCM?
    \end{itemize}
  }
\end{frame}


%--------------------------------------------------------------------------------------------------------------
\begin{frame}[t]
  \frametitle{Outline}

  \tableofcontents
  \vfill
  \begin{flushright}
    \tiny{\textbf{General Reference:} \cite{jacobs2010operations}~\cite{ballou2007evolution} \cite{chopra2007supply} \cite{goetschalckx2011supply} } \hskip3ex
  \end{flushright}

  %%%------
  \note{
    \textbf{Purpose:} overall view of lecture\\
    \textbf{Leading Question:} $<<None>$ \\
    \begin{itemize}
      \item \textbf{Understand:} limit man/time/computation power $\rightarrow$ what should we focus
      \item \textbf{Forecast:} give me some number to work with $\rightarrow$ art\&science of systematic guessing demands
      \item \textbf{Oper.Plan:} how to smooth production $\rightarrow$ overall concept of planning
      \item \textbf{Aggregate:} overall picture of family demand $\rightarrow$ what is your strategy (demand chase/inventory/outsouces)
      \item \textbf{S\&OP:} integrate Capacity and department involved in planning and try to influence demands
    \end{itemize}
  }
\end{frame}

%%=================================================================================================
\begin{frame}
  \frametitle{Game Plan: From Conception to Shipping...}


  \begin{itemize}
    \item Identifying \& assessing our \alert{capability} $\Longleftrightarrow$ Strategic planning \& mgt
    \item Guessing customer's \alert{demands} \& \alert{requirements} $\Longleftrightarrow$ Forecasting
    \item Smoothing \& Managing \alert{resources} $\Longleftrightarrow$ Aggregate planning \& Demand mgt
    \item \alert{Sourcing} R/M needed to produce $\Longleftrightarrow$ Purchasing \& Material requirement
    \item<2- |handout:1> Estimating \alert{producing time} $\Longleftrightarrow$ Production planing \& Manuf mgt
    \item<2- |handout:1> Managing \alert{stocks}  $\Longleftrightarrow$ Warehousing mgt \& Inventory control
    \item<2- |handout:1> Maintaining availability \& service levels $\Longleftrightarrow$ Channel \& Distribution mgt
    \item<2- |handout:1> Ensuring commitment \& \alert{delivery} products $\Longleftrightarrow$ Routing, Outsoucing, \& Shipping
  \end{itemize}

  %%%------
  \note{
    \textbf{Purpose:} link practical questions with discipline\\
    \textbf{Leading Question:} which topics in textbook for practical questions\\
    \begin{itemize}
      \item $<<Explain Each Bullet>>$
      \item In this lecture, we will focus on forecasting and demand management
    \end{itemize}
    Other areas that may Less familiar with IE
    \begin{itemize}
      \item Capturing niche market $\Longleftrightarrow$ Product \& process design
      \item Maintaining customer loyalty $\Longleftrightarrow$ Quality
      \item Having utilization $\Longleftrightarrow$ Maintenance
    \end{itemize}
  }
\end{frame}

%================================================================================================-
\section[Demands101]{Understanding your demands}
%================================================================================================-
\begin{frame}
  \frametitle{What are demands? }

  \begin{alertblock}{Are these terms difference?}{}{}
    \begin{center}
      demand $=\!\!\!\!?$ sale $=\!\!\!\!?$ order $=\!\!\!\!?$ forecast \\
    \end{center}
  \end{alertblock}

  \NextLine \NextLine
  \textsc{Nature of demands}
  \begin{itemize}
    \item \textbf{Price Elasticity:} price dependence (electricity, TH cigarette)
    \item \textbf{Predictability:} fluctuation, seasonal, customer-driven (meme)
    \item \textbf{Service Lv:} expectation (iPhone), customization
    \item \textbf{Geographic:} regional, location
  \end{itemize}

  %%%------
  \note{
    \textbf{Purpose:} clarify terms and why we always focus on demand\\
    \textbf{Leading Question:} How these term are difference? What and When you know? \\
    \begin{itemize}
      \item \textbf{Demand:} qty in customers' mind that farecast tries to capture. some may does not know that it needed.
      \item \textbf{Sales:} qty that customers actually purchase and record. it usually less than demand
      \item \textbf{Order:} qty that customers order and lead to sales. some company assume that equivalent to demand
      \item \textbf{Forecast:} qty that company tries to guess based on sales and which to be the same as order/demand.
    \end{itemize}
    \textsc{Thinking Points:}
    \begin{itemize}
      \item \textbf{Resources:} clothes, ticket, passport, money, medication, shoes, cosmetic
      \item \textbf{Factor:} day, location, purpose, whom travel,
      \item \textbf{Manage:} money and time

    \end{itemize}
  }
\end{frame}

%================================================================================================-
\begin{frame}
  \frametitle{How should we manage demands?}

  \only<1 |handout:1>{
  \begin{center}
    \includegraphics[width=0.50\linewidth]{\figureLocation{scm_demand/demandPlanning.png}}\\
  \end{center}
  }

  \only<2-|handout:0>{
  \begin{itemize}
    \item Understanding \alert{nature of demand}
    \item Prioritizing \alert{independent easy to monitor} demand %Should all demands treated equally? % prioritizing demand
    \item Forecasting \alert{`reasonable'} demand %ing Since ${demand} = Pr(\cdot)$, should we ignore planning? % planning demand
    \item \alert{Communicating and updating} demand forecast % communicate demand
    \item \alert{Responding} to expectation of demands % communicate demand
    \item \alert{Influencing} quantity and time of demand %Can we influence demands? If so, how? % p
  \end{itemize}
  }

  %%%------
  \note{
    \textbf{Purpose:} concept and overview of demand management \\
    \textbf{Leading Question:} How should we manage demands?\\
    \begin{itemize}
      \item many view SCM as demand chain management (how we translate and communicate consumer's demand to every stakeholder)
      \item demand prediction = holy grail in business
      \item \textbf{\textsl{IE/OM Concept:}}focus things that matter, estimate the rest
      \item \textbf{Nature:} an occupation in IE require understand and exploit pattern. Never treat everything equal $\rightarrow$ affect more with less
      \item \textbf{Independent:} less should have independent and important $\rightarrow$ paerto diagram of some items with certain features
      \item \textbf{Forecast:} use data (itself and relevant factors) to predict future demand
      \item \textbf{Communicate:} $\exists$ many uncertainty $\rightarrow$ feedback often with simple forecast model
      \item \textbf{Respond:} ; second market
      \item \textbf{Influence:} during the process we may can influence $\rightarrow$ balance demand. Area of \textsl{Revenue Management}
    \end{itemize}


  }
\end{frame}

%================================================================================================-
\begin{frame}
  \frametitle{Understanding nature of demands:}

  \begin{itemize}
    \item Aggregate VS \textbf{Market}: Why do we think about demand of \alert{our brand}?
    \item Stochastic VS  \textbf{Deterministic}: Why do we think about \alert{constant} demand?
    \item Dependent VS \textbf{Independent}: Why do we think about \alert{independent} demand?
  \end{itemize}

  \NextLine
  \textsc{How do we manage/model inventory?  }
  \begin{center}\small{
  \begin{alignat*}{4}
    \mbox{Demand}^{dep}~~=& ~f(\mbox{Demand}^{ind}_i,\mbox{Demand}^{ind}_j,\cdots) \\
    \mbox{Demand}^{ind}_i~~=&  ~\mbox{demand}_{certain} ~~ &~+~&  \mbox{demand}_{random} \\
                    &  ~~\Downarrow             ~~&~~~&  ~~ \Downarrow \\
                    &  \bullet steady          ~~ &~~~&  \bullet safety~stock \\
                    &  \bullet trend \\
                    &  \bullet seasonal
  \end{alignat*}}
  \end{center}

  \begin{itemize}
    \item[] \textbf{Note:} consolidate demands $\rightarrow$ better mgt and forecasting  \hfill \hyperlink{postpone}{\beamerbutton{Pg:HP Postponement}}
  \end{itemize}

  %%%------
  \note{
    \textbf{Purpose:} overall view of IT in logistics\\
    \textbf{Leading Question:} How each IT system related to one other?\\
    \begin{itemize}
      \item \textbf{Market:} only demand of our product in certain channel, not total (\textsl{Aggregate})
      \item \textbf{Deterministic:} forecast an average or very likely demand, not range of demand (\textsl{stochastic})
      \item \textbf{Independent:} focus only demand of independence items matter, not every item (\textsl{dependent})
      \item \textbf{Component of Demands:}
      \begin{itemize} \tiny{
        \item \textbf{\textsl{centain:}} steady (fixed level), trend (upward/downward), seasonal (up-down with advance knowledge), economic cycle (long term seasonal)
        \item \textbf{\textsl{random:}} safety stock (excess for uncertainty), }
      \end{itemize}
    \end{itemize}
  }
\end{frame}

%================================================================================================-
\begin{frame}[label=postpone]
  \frametitle{Postponement: delay differentiation \forIE}

  \textsc{Motivation:} HP Customization
  \begin{itemize}
    \item \textbf{Motivation:} inventory, nature of demand, multiple markets
    \item \textbf{Idea:} \alert{delay the differentiation} of product/process
    \item \textbf{Related concept:} localization, simple and modular design
    \item \textbf{Trade-offs:} economy of scale, production cost \textbf{VS} inventory, transportation
    \item \textbf{Implementation:} cross-organization coordination (manufacturing, purchasing, R\&D, market)
    \begin{itemize}
      \item \textbf{\textsl{Product design:}} reducing inventory \& leadtime by std parts \& manuf.
      \item \textbf{\textsl{Process design:}} delaying process differentiation by technology
      \item \textbf{\textsl{Agile supply chain:}} time based competitive
    \end{itemize}
    \item \textbf{Example:} DeskJet, Paint mixture, Benetton shirt, HP hard drive test
  \end{itemize}

  \note{
    \textbf{Purpose:} a classical paper about forecasting and postponement at HP, this paper sparks \alert{agile SCM}, \alert{postponement}, \alert{std testing}
    \begin{itemize}
      \item \textbf{purpose:} re-cap experience at HP to re-think how inventory and demand management should work
      \item \textbf{re-think about inventory:} what form, nature, and channel-market that it serves
      \item \textbf{idea:} de-couple product-process. ask why?
      \item \textbf{localization:} ask local vendor to produce bulky simple component to reduce transportation and risk
      \item \textbf{re-design:} reduce component and possibility of deadstock by redesign and family group
      \item \textbf{Independent modular product design}
      \begin{itemize} \tiny{
        \item maximizing standard parts $\rightarrow$ reduce inventory
        \item manufacturing simultaneously $\rightarrow$ reduce leadtime }
      \end{itemize}
      \item \textbf{Independent modular process design}
      \begin{itemize} \tiny{
        \item \textbf{paint mixing machine:} color analyzing machine + base color mixing machine
        \item \textbf{benetton dyeing:} delay dyeing process until realizing customer demands
        \item \textbf{HP hard drive testing process:} standard testing and supplier testing }
      \end{itemize}
     \item \textbf{Agile supply chain}
      \begin{itemize} \tiny{
        \item hiring local for value added process (e.g., packaging, assembly, labeling)
        \item reducing inventory and its permutation by postpone assembly process
        \item stock parts and local capacity }
      \end{itemize}
    \end{itemize}
    }
\end{frame}

%================================================================================================-
\begin{frame}
  \frametitle{Prioritizing demand with parieto principle}

  \begin{itemize}
    \item \textbf{Idea:} \alert{few items} affects \alert{majority of activity}
    \item[] ~
    \begin{itemize}
      \item[] \hskip5ex \small{\emph{You don’t get what you want in life}}
      \item[] \hskip3ex \small{\emph{You get what you focus on and expect.}}
      \item[] ~\hskip18em \tiny{---Tony Rush} ~\\[8pt]
    \end{itemize}
    \item \textbf{Therefore:} focus on \alert{few} that are \alert{important}
    \item \textbf{Application:} ABC analysis
    \begin{description}
      \item[Class A] important (15-20 percentile $\rightarrow$ for 55-75\% of activity)
      \item[Class B] some-what important (20-50 percentile $\rightarrow$ 15-15\% of activity)
      \item[Class C] not important (50-100 percentile $\rightarrow$ 5-15\% of activity)
    \end{description}
  \end{itemize}

  \note{
    \textbf{Purpose:} re-cap parieto principle\\
    \textbf{Leading Question:} how to measure which items are matter\\
    \begin{itemize}
      \item \textbf{Parieto's Observation} only 20\% affects 80\% movement, cost
      \item \textbf{How to measure:} value = qty $\times$ unit price
      \item \textbf{Warning:} some application only qty (AS/RS in warehouse) or price matter (project)
    \end{itemize}

  }
\end{frame}

%================================================================================================-
\begin{frame}
  \frametitle{Example of ABC analysis}

  \textsc{}
  \begin{center}
  \footnotesize{
  \begin{tabular}{lrr}\hline
  Item No.  & Annual dollar   & \% Total Value \\ \hline
  22        & 95,000          & 40.69 \\
  68        & 75,000          & 32.13 \\
  27        & 25,000          & 10.71 \\
  03        & 15,000          & 6.43 \\
  82        & 13,000          & 5.57 \\
  54        & 7,500           & 3.21 \\
  36        & 1,500           & 0.64 \\
  19        & 800             & 0.34 \\
  23        & 425             & 0.18 \\
  41        & 225             & 0.10 \\ \hline
  -         & \$223,450       & 100\% \\ \hline
  \end{tabular}
  }
  \end{center}

  \begin{center}
    \includegraphics[width=0.65\linewidth]{\figureLocation{scm_demand/ABC.jpg}}\\
  \end{center}
  \begin{flushright}
    \tiny{\textbf{source:} Chase and Jacob. 2011.\hskip10em}\\
  \end{flushright}

  \note{
    \textbf{Purpose:} example of calculation  \\
    \textbf{Leading Question:} $<<$None$>>$\\
    \begin{itemize}
      \item SCM
    \end{itemize}
    \begin{itemize}
      \item A: 22, 68   $\rightarrow$ 72.9\%
      \item B: 27, 03, 82 $\rightarrow$ 22.7\%
      \item C: 54, 36, 19, 23, 41 $\rightarrow$ 4.4\%
    \end{itemize}
  }
\end{frame}

%================================================================================================-
\section[Forecasting]{Forecasting SCM Demand}
%================================================================================================-
\begin{frame}
  \frametitle{What should know on forecasting }

  \begin{itemize}
    \item \textbf{What:} A process to understand \alert{systematic of demands}
    \item \textbf{Facts:} Inaccurateness, but all companies need it
    \begin{itemize}
      \item Forecasting is, typically, \alert{incorrect}
      \item Forecasting is suitable for a \alert{group of products}
      \item Forecasting is \alert{inaccurate as time horizon increases} \NextHalfLine
    \end{itemize}
    \item \textbf{Ideal:} Valid for \alert{short time} and \alert{need updated}
    \item \textbf{Aware} Art \& Science: roles of marketing \& sale; common misconcept
    \item \textbf{Types:} Qualitative, Time Series, Regression, Simulation
  \end{itemize}

  \NextLine
  \begin{block}{\small{\textsc{Need Help!}}  }{}
    \begin{itemize}
      \item \textbf{Recap:} \href{https://ie.eng.chula.ac.th/~oran/classes/LSCM/presentation/scm_forecasting.pdf}{\code{04 time series (extra)}}
      \item \textbf{Textbook:} \href{https://otexts.com/fpp3/}{\code{https://otexts.com/fpp3/}}
    \end{itemize}
  \end{block}


  \note{
    \textbf{Purpose:} re-cap forecasting concept, esp rules and warning\\
    \textbf{Leading Question:} What do you know about forecasting\\
    \begin{itemize}
      \item \textbf{systematic} traceable/repeatablity/transpalentcy/explanatory
      \item \textbf{concept:} guess future by hoping that it follow same pattern and randomness will cancel each other
      \item \textbf{Misconcept}
      \begin{itemize} \tiny{
        \item let \alert{model/program} run
        \item adjust forecasting based on supply/ capacity
        \item wrong type of forecasting, e.g., operation \textsc{VS} budget, sku \textsc{VS} family }
      \end{itemize}
      \item \textbf{Qualitative} use experts to deliberate on qty of demand, e.g. Delphi's method, sales opinions (incentive), Consumer surveys (new product), Outside opinion
      \item \textbf{Time Series} use historical qty to predict itself $\rightarrow$ high correlation between historical data. Major part of this lecture, e.g. stock price
      \item \textbf{Regression} use other factors to predict $\rightarrow$ relation with factor, good for explanatory, e.g., using GDP and crude oil to predict consumption of electricity
      \item \textbf{Simulation} combine both time series and regression to predict, e.g. weather forecast using moister, wind, heat combine with previous temperature to predict
    \end{itemize}
  }
\end{frame}

%%============================================================================================================
\begin{frame}
  \frametitle{What should aware before forecast }

  \begin{itemize}
    \item \textbf{Description:} story, relationship with other data
    \item \textbf{Time Horizon:} hour, day, week, year
    \item[] \hskip5ex Actual = Forecast + error
    \item \textbf{Pattern of Data:} seasonal, trend, cycle
    \item \textbf{Forecasting Model:} \alert{assumptions}, data required, parameters (static \textsc{VS} dynamic)
    \item \textbf{Accuracy:} measuring, how to improve
  \end{itemize}

  \note{
    \textbf{Purpose:} overall view of IT in logistics\\
    \textbf{Leading Question:} How each IT system related to one other?\\
    \begin{itemize}
      \item \textbf{What are forecasting for}
      \begin{itemize} \tiny{
        \item \textbf{Workforce and Capacity:} hiring, training
        \item \textbf{Operation:} purchasing, scheduling, WIP, negotiation, planning
        \item \textbf{Accounting:} cash flow, equipment investment
        \item \textbf{Marketing:} pricing, promotion, place }
      \end{itemize}
      \item Forecasting $\neq$ Math; correct number/formulation; accuracy
      \item \textbf{Description:} story makes forecasting real and alive
      \item \textbf{Time Horizon:} how often to make prediction $\rightarrow$ how often to make decision and alignment of data
      \item \textbf{Pattern:} discover by plotting
      \item \textbf{Model:} what are assumptions and data required? Underlying pattern. how parameters interact and "memories"
    \end{itemize}
  }
  \end{frame}

%%================================================================================================-
\begin{frame}
  \frametitle{Example: US Air Passengers 1949-1951 }

  \RemoveLine
  \only<1|handout:1>{
  \begin{center}
    \includegraphics [width=\textwidth]{\figureLocation{scm_demand/AirPassengers.pdf}}
  \end{center}
  }
  \only<2|handout:0>{
    \begin{center}
      \includegraphics [width=\textwidth]{\figureLocation{scm_demand/AirPassengersYearly.pdf}}
    \end{center}
  }

  \note{
    \begin{center}
      \includegraphics [width=0.3\textwidth]{\figureLocation{scm_demand/AirPassengersYearly.pdf}}
    \end{center}
    \textbf{Purpose:} remind students about pattern and importance of story\\
    \textbf{Leading Question:} What do you see in this slide?\\
    \begin{itemize}
       \item what do we want to forecast? any data? guess?
       \item point out smooth, trend, and seasonality
       \item 1950 Airline Industry was niche market; After WWII
       \item Advent of Jet Aircraft, Summer and Winter travel
    \end{itemize}
  }
\end{frame}

%============================================================================================================--
%%% Notation %%%
\newcommand{\Forecast}{F}
\newcommand{\Actual}{A}
\newcommand{\Error}{E}
\newcommand{\Smooth}{S}
\newcommand{\Index}{r}
\newcommand{\Data}{X}
%%================================================================================================-
\begin{frame}
  \frametitle{Time Series Review: Workshop }

  \begin{itemize}
    \item \textbf{Forecasting Accuracy:} Actual$_t$ - Forecast$_t$ (i.e., MAPE, MSE)
    \item \textbf{Naive Method:}  Forecast \alert{=} actual
    \item \textbf{Smoothing Tech:} Forecast \alert{$\approx$} set of actual
    \begin{itemize}
      \item \textbf{Moving Average:} Forecast = \alert{average} of actual
      \item \textbf{Exponential Smoothing} Forecast = \alert{combination} of actual and \alert{last forecast}
    \end{itemize}
    \item \textbf{Trending:} Forecast = expo smoothing +  \alert{slope} $\times$ period
    \item[] \hskip10ex~     Slope = \alert{gap} of forecast
    \item \textbf{Seasonality:} Forecast$_{sea}$ = $\Index_{sea} \cdot $ Forecast$_{de-sea}$
  \end{itemize}


  \note{
    \textbf{Purpose:} summary concept of all importance forecasting model\\
    \textbf{Leading Question:} $<<$None$>>$\\
    \begin{itemize}
      \item \textbf{Accuracy:} always respect actual. MAPE $= \frac{|A_t - F_t|}{A_t}$ and MSE $= (A_t-F_t)^2$
      \item \textbf{Naive:} use when few data
      \item \textbf{Smoothing:} use when no pattern data or try to remove noise
      \item \textbf{Trend:} use when $\exists$ trends. also call double exponential smoothing
      \item \textbf{Seasonal:} use when $\exists$ seasonal pattern more than two cycles with the idea of remove effect of seasonality before do forecasting
    \end{itemize}

  }
\end{frame}

%================================================================================================-
\begin{frame}
  \frametitle{Communicating and updating demand forecast}

  \begin{quotation}
    If you have to forecast, forecast often.
  \end{quotation}
  \begin{flushright}
    \hfill \tiny{\textbf{source:} Edgar Fiedler \hskip10ex~}
  \end{flushright}

  \NextHalfLine
  \textsc{Sharing Points:}
  \begin{itemize}
    \item \textbf{Know data source:} market goal $=\!\!\!\!?$ operation cap, reliability \textsc{VS} urgency
    \item \textbf{Keep process simple/ comm. flat:} verbal \& constant, formal \& informal, time frame
    \item \textbf{Right communication channel:} always confirm \& name (scope), IT as enable (goldcity)
  \end{itemize}

  \begin{block}{}{}
    Believe the forecasting pattern, not absolute \alert{quantity} or \alert{time}
  \end{block}

  \note{
    \textbf{Purpose:} remind about communication as we apply a simple forecasting model\\
    \textbf{Leading Question:} How to handel inaccuracy and uncertainty during forecasting?\\
    \begin{itemize}
      \item \textbf{Source:} must know who gives forecasting and their agenda as well as reliability of sources
      \item \textbf{How to manage urgency:} impact and alternative of vendors
      \item \textbf{Simple Process:} select formal and informal by time frame; use verbal comm and keep every one in the loop
      \item \textbf{Comm Channel:} remind that comm is two-way process always confirm name and contact info; goldcity uses IT as a formal way to communicate sale order and verifies credit
      \item \textbf{Operation Cap:} \# OT, storage space, machine cap, yield, loading
    \end{itemize}
  }
\end{frame}

%================================================================================================-
\begin{frame}
  \frametitle{Strategies for Responding to demands}

  \begin{quotation}
     Supply always comes on the heels of demand.
  \end{quotation}
  \begin{flushright}
    \hfill \tiny{\textbf{source:} Robert Collier \hskip10ex~}
  \end{flushright}

  \NextHalfLine
  \textsc{Sharing Points:}
  \begin{itemize}
    \item \textbf{Agile Design:} modular product, delayed differentiation, flexible capacities
    \item \textbf{Responsive Sourcing:} diversification, strategic partnership
    \item \textbf{Adaptive Network:} omni-channel, dyn routing, collaborative network
    \item \textbf{Data-Driven Sensing:} adv forecasting tech, scarcity \& halo technique
  \end{itemize}


  \note{
    \textbf{Purpose:} remind about communication as we apply a simple forecasting model\\
    \textbf{Leading Question:} How to handel inaccuracy and uncertainty during forecasting?\\
    \begin{itemize}
      \item \textbf{Source:} must know who gives forecasting and their agenda as well as reliability of sources
      \item \textbf{How to manage urgency:} impact and alternative of vendors
      \item \textbf{Simple Process:} select formal and informal by time frame; use verbal comm and keep every one in the loop
      \item \textbf{Comm Channel:} remind that comm is two-way process always confirm name and contact info; goldcity uses IT as a formal way to communicate sale order and verifies credit
      \item \textbf{Operation Cap:} \# OT, storage space, machine cap, yield, loading
    \end{itemize}
  }
\end{frame}

%================================================================================================-
\begin{frame}
  \frametitle{Influencing demand}

  \begin{quotation}
    Cutting prices or putting things on sale is not sustainable business strategy.
  \end{quotation}
  \begin{flushright}
    \hfill \tiny{\textbf{source:} Howard Schultz \hskip10ex~}
  \end{flushright}


  \textsc{Sharing Points:}
  \begin{itemize}
    \item \textbf{Other factors:} service  (physician schedule), experience (food), waiting time (BTS), events (MOTOR SHOWS), quality (GRAB)
    \item \textbf{Marketing:} focused social media (time spent)
  \end{itemize}

  \note{
    \textbf{Purpose:} overall view of IT in logistics\\
    \textbf{Leading Question:} How each IT system related to one other?\\
    \begin{itemize}
      \item SCM
    \end{itemize}
  }
\end{frame}


%==============================================================================================================
\section[Planning]{Operational Planning Concepts}
%============================================================================================================--
\begin{frame}
  \frametitle{What is planning?}

  \begin{description}
    \item[Planning:] Process of determining \alert{how} an organization can get \alert{destination} it wants to go and what \alert{actions} it will do to accomplish its objectives. Through planning, the organization identifies:
      \begin{itemize}
        \item Where it/we is going?
        \item How it/we will get there?
      \end{itemize}
  \end{description}

  \NextLine
  \textsc{Thinking Points:} Recall your last travelling experience
  \begin{itemize}
    \item<2-> What do you need (\alert{resources}) ?
    \item<2-> How to \alert{estimate} such needed resources?
    \item<2-> How you \alert{manage} if you estimation is incorrect?
  \end{itemize}




  %%%-------------------
  \note{
    \textbf{Purpose:}  \\
    \textbf{Leading Question:} \\
    \textsc{Purposes of Planning}\\
    \begin{itemize}
      \item Minimize risk by \alert{spell-out} uncertainties $\rightarrow$ future-oriented, predict \alert{outcomes}
      \item Increase organizational success $\rightarrow$ \alert{clarify} objectives, procedure, possible problems
      \item Establish \alert{coordinated effort} $\rightarrow$ decision \alert{coordination}
    \end{itemize}
    }
\end{frame}

%============================================================================================================--
\begin{frame}
  \frametitle{Operation Planning Concept}

  \begin{itemize}
    \item \textbf{Hierarchy:} decision in upper-level $\rightarrow$ constraint in lower-level
    \item \textbf{Feedback:} systematic correction from its result
    \item \textbf{'Single' database \& Integration:} maintaining data integrity and accuracy
    \item \textbf{Coverage:} cover 80\% of cases or focus on \alert{bottleneck}
    \item \textbf{Flexibility:} \alert{alternative}, answering what if, provide \alert{exception}
    \item \textbf{Transparency:} understanding logic and algorithm behind system
    \item[]
    \item<2-> \textbf{OKC's Note:} keep it simple and roll quickly
  \end{itemize}

  \NextHalfLine
  \uncover<3->{
  \begin{block}{\small{Warning!!}}{}
    \begin{center}
    Be positive; Everyone makes mistake
    \end{center}
  \end{block}
  }

  \note{
  \begin{center}
    \includegraphics [height=0.30\textheight]{\figureLocation{scm_demand/planningRel.jpg}}
  \end{center}
    \textbf{Purpose:} re-cap planning concept in IE/OM\\
    \textbf{Leading Question:} What is a good plan look like\\
  \begin{itemize}
    \item \textbf{Hierarchical:} from big picture planning (Aggre.Plan) to detail planning (Oper.Plan); detail needed when almost execute
    \item \textbf{Feedback:} use capacity and material as feed back to change within level plan, then upper/lower level plan. no direct capacity or material interaction
    \item \textbf{Man-Machine:} expect no perfect plan; leave area for exception, esp human nature
    \item \textbf{Database:} all parameters and planning constants should be available in a single accurate database
    \item \textbf{Transparency:} everyone should be able to repeat and duplicate plan given same information
    \item \textbf{Aware:} Planning as Communication tools; Planning $\approx$ Forecasting $\rightarrow$ incorrect/guesssing by nature
    \item \textbf{Specific response time:} designing for improvement
  \end{itemize}
  }
\end{frame}

%==============================================================================================================
\begin{frame}
  \frametitle{Production Process}

  \begin{center}
    \includegraphics [width=0.65\textwidth]{\figureLocation{scm_demand/productProcessMatrix.jpg}}
  \end{center}
  \begin{flushright}
    \tiny{\textbf{source:} Chase \& Jacabs. 2010. \emph{Operation and Supply Chain Management}. } \hskip10em
  \end{flushright}


  \scriptsize{
  \begin{itemize}
    \item \textbf{Project:} product remains in a fixed location \uncover<2->{ \hfill \code{Gantt Chart}}
    \item \textbf{Work center:} similar equipment or functions are grouped together  \uncover<2->{ \hfill \code{Machine Scheduling}}
    \item \textbf{Manufacturing cell:} dedicated area where similar products are produced
    \item \textbf{Assemble line:} processes are arranged according to the progressive steps
    \item \textbf{Continuous process:} assembly line only the flow is continuous \uncover<2->{ \hfill \code{Model Change-Over}}
  \end{itemize}
}

  \note{
    \textbf{Purpose:} re-cap production process matrix and how to apply planning concept\\
    \textbf{Leading Question:} What area of production that the planning in this class focus?\\
    \begin{itemize}
      \item \textbf{Explanation:} for forgetful IE student. This is the re-cap of production process matrix. What suitable process given a characteristics of a product.
        \begin{itemize} \tiny{
          \item \textbf{project:} ETO or MTS, e.g., air plane, large building, industry machine
          \item \textbf{Work Center (flow shop):} generalize machines for multi-purpose (each machine require a single skilled worker), e.g., metal lathe machine, job shop
          \item \textbf{Manuf cell:} grouped in machine connected together into cell (multiple machines operates by single worker)
          \item \textbf{Assemble line:} a mass production with conveyor for a family of product, e.g., car assembly
          \item \textbf{continuous process:} a mass production dedicated machine with no variation and product mixed, e.g., bottle water,  }
        \end{itemize}
      \item \textbf{Mass Customization:} for high volume with individuate specification, e.g., contact lens, paint mixing, custom fip-flop shoe
      \item \textbf{3D Printing} for low volume with high standard,
      \item \textbf{Our Planning Focus:} Work Center, Manufacturing Cell, and Assemble Line
    \end{itemize}


  }
  \end{frame}

%%================================================================================================-
%% Standardize how tikz displays vendors, crossdock, and stores
\usetikzlibrary{arrows,positioning,decorations.pathmorphing,calc,shapes}
%%============================================================================================================
%\begin{frame}
%  \frametitle{Time Horizon}
%
%  \begin{center}
%  \begin{tikzpicture}[xscale=0.75,yscale=0.75,>=stealth]
%    \tikzstyle{Lecblock}=[rectangle,draw=black!80,thick,fill=red!15,text centered, text width=28mm,text height=7mm,align=left];
%
%    \node[Lecblock] (Inter)                    {\textFont{
%      \begin{tabular}{l}
%        $\bullet$ Forecasting \\
%        $\bullet$ Manpower \\
%        $\bullet$ Distribution \\
%        ~~
%      \end{tabular} } };
%    \node[Lecblock] (Short) [left= of Inter]   {\textFont{
%      \begin{tabular}{l}
%        $\bullet$ Purchasing \\
%        $\bullet$ Scheduling \\
%        $\bullet$ Maintenance \\
%        $\bullet$ Inventory
%      \end{tabular} } };
%    \node[Lecblock] (Long)  [right= of Inter]   {\textFont{
%      \begin{tabular}{l}
%        $\bullet$ Facility \\
%        $\bullet$ Sale channels \\
%        $\bullet$ Service \\
%        $\bullet$ Objective
%      \end{tabular} } };
%
%    \node[] (tInter) [above= of Inter,yshift=-10mm] {\textFont{\textbf{Intermediate}}  };
%    \node[] (tShort) [above= of Short,yshift=-10mm] {\textFont{\textbf{Short-term}} };
%    \node[] (tLong)  [above= of Long, yshift=-10mm]  {\textFont{\textbf{Long-term}} };
%    \node[] (tHorizon) [above= of Inter] {\small{\textbf{Decision horizon}} };
%    \node[] (tIssue)   [below= of Inter] {\small{\textbf{Issues}} };
%
%    \draw[->,thick] (tHorizon.west) -- (tShort.north);
%    \draw[->,thick] (tHorizon.south) --(tInter.north);
%    \draw[->,thick] (tHorizon.east) -- (tLong.north);
%    \draw[->,thick] (tIssue.west) -- (Short.south);
%    \draw[->,thick] (tIssue.north) --(Inter.south);
%    \draw[->,thick] (tIssue.east) -- (Long.south);
%  \end{tikzpicture}
%  \end{center}
%
%  \begin{flushright}
%    \tiny{\textbf{adopted from:} \textsl{Nahmias, S.} 2001.}
%  \end{flushright}
%
%  \note{
%    \begin{itemize}
%      \item \textbf{Planning horizon:} difficult to change?
%      \item \textbf{long term:} quantities \& instince
%      \item \textbf{short term:} relay data \& solving day-to-day problems
%      \item \textbf{intermediate:} optimize
%    \end{itemize}
%  }
%\end{frame}


%============================================================================================================
\begin{frame}
  \frametitle{Closed-Looped Production Planning}

  \begin{center}
  \begin{tikzpicture}[xscale=0.75,yscale=0.75,>=stealth]

  \tikzstyle{FeedBack}=[rectangle,draw=black!80,thick,fill=green!15,text centered,text width=18mm,text height=2mm];
  \tikzstyle{Planning}=[rectangle,draw=black!80,thick,fill=blue!15,text centered,text width=18mm,text height=2mm];
  \tikzstyle{Info}=[rectangle,draw=black!80,thick,fill=red!15,text centered,text width=18mm,text height=2mm,rounded  corners=1ex];
  \tikzstyle{Output}=[rectangle,draw=black!80,thick,fill=orange!15,text centered,text width=18mm,text height=2mm];


    \node[Planning] (MS) {\textFont{MPS}};
    \node[Info] (DemandMgt)  [left= of MS] {\textFont{Demand Mgt.}};
    \node[Info] (CustOrder)  [left= of DemandMgt, yshift=8mm ,xshift=25mm] {\textFont{PO}};
    \node[Info] (ShortFore)  [left= of DemandMgt, yshift=-8mm,xshift=25mm] {\textFont{S. Forecasting}};
    \node[FeedBack] (RCCP) [right= of MS] {\textFont{RCCP}};


    \node[Planning] (PP)  [above= of MS, yshift=8mm] {\textFont{Agg. Planning}};
    \node[Info] (LongFore) [left= of PP]  {\textFont{L. Forecasting}};
    \node[FeedBack] (RRP) [right= of PP]  {\textFont{RRP}};
    \coordinate (top) at ($(CustOrder.north)+(-1,0.5)$) ;


    \node[Planning] (MRP) [below= of MS, yshift=-8mm] {\textFont{MRP}};
    \node[Info] (BOM) [left= of MRP, yshift=4mm] {\textFont{BOM }};
    \node[Info] (Inv) [left= of MRP, yshift=-4mm] {\textFont{Inventory}};
    \node[FeedBack] (CRP) [right= of MRP] {\textFont{CRP}};
    \coordinate (bottom) at ($(ShortFore.south)+(-1,-0.5)$) ;

    \node[Output] (Job) [below= of MRP,xshift=-15mm] {\textFont{Job Shop}};
    \node[Output] (Pur) [below= of MRP,xshift=15mm] {\textFont{Purchasing}};

    %vertical
    \draw[<->,thick] (MS) -- (PP);
    \draw[<->,thick] (MS) -- (MRP);
    \draw[<->,thick] (Job) -- (MRP);
    \draw[<-,thick]  (Pur) -- (MRP);

    %Lv1 horizontal
    \draw[<-,thick] (PP) -- (RRP);
    \draw[<-,thick] (PP) -- (LongFore);

    %Lv2 horizontal
    \draw[<->,thick] (MS) -- (RCCP);
    \draw[<->,thick] (DemandMgt) -- (MS);
    \draw[<-,thick] (DemandMgt) -- (ShortFore);
    \draw[<-,thick] (DemandMgt) -- (CustOrder);

    %Lv3 horizontal
    \draw[<->,thick] (MRP) -- (CRP);
    \draw[<-,thick] (MRP) -- (BOM);
    \draw[<-,thick] (MRP) -- (Inv);

    %Separation
    \draw[black!20,dashed]  (top)   --node[black!30,above, pos=0.3]{\small{Long-Term Plan}} ++(13,0) ;
    \draw[black!20,dashed]  (bottom) --node[black!30,above, pos=0.3]{\small{Intermediate Plan}} ++(13,0) ;
    %\node[black!30] () [below=of top, xshift=6mm] {\small{Shot-Term Plan}};

%    \draw[->,thick]  (DRP.south east) -- (MS.north west);
  \end{tikzpicture}
  \end{center}
  \begin{flushright}
    \tiny{\textbf{adopted from:}  \textsl{Smith, S.} 1989.}
  \end{flushright}


  \note{
    \textbf{Purpose:} overall view of IT in logistics\\
    \textbf{Leading Question:} How each IT system related to one other?\\
    \begin{itemize}
      \item \textbf{Re-Cap:} for forgetful IE/OM students. Here is each planning mean
  \begin{itemize} \tiny{
    \item \textbf{Aggregate Planning:} Family item planning for key resources or manuf policies
    \item \textbf{Bill Of Material (BOM):} show lead time \& qty.
    \item \textbf{Master Production Scheduling (MPS):} FG planning
    \item \textbf{Material Requirement Planning (MRP):} All materials planning with leadtime
    \item \textbf{Resource Requirements Planning (RRP):} system no barriers
    \item \textbf{Rough Cut Capacity Planning (RCCP):} use average set up time, with Bill of capacity (BOC)
    \item \textbf{Capacity Requirements Planning (CRP):} difficult to use  }
    \end{itemize}

  \end{itemize}
  }
  \end{frame}

%============================================================================================================
\section[Aggre. planning]{Aggregate Planning }
%============================================================================================================

\begin{frame}
  \frametitle{Aggregate Planning}

    \begin{itemize}
      \item \textbf{What:} planning for \alert{family of product}
      \item \textbf{Idea:} smoothing demands \& key resources (workers, \# machine)
      \item \textbf{Occurred:} beginning of fiscal year (\emph{unit:} month of sale)
      \item \textbf{Input:} long-term forecasting + key resources, policy (overtime, subcontract, backlog, machine capacity)
      \item \textbf{Output:} production plan, capacity, inventory at \alert{acceptable cost}
      \item \textbf{Note:} must be flexible and share-able
    \end{itemize}

  \note{
    \textbf{Purpose:} overview of aggregate planning \\
    \textbf{Leading Question:} $<<$None$>>$ \\
    \begin{itemize}
      \item \textbf{What for:} family of product best for 3-18 months
      \item \textbf{why important:} every party/company needs a road map, put information to good use
      \item \textbf{Options:} capacity, utilization, inventory
      \item \textbf{Note:} Many uses Optimization for Aggregate planning
      \item \textbf{Trials:} SOP \emph{manipulate demands \& supplies}
    \end{itemize}
  }
\end{frame}

%============================================================================================================
\begin{frame}
  \frametitle{Red Tomato: Aggregate planning}

  \small{
    Red Tomato Tool, a US gardening tool manufacturer, usually faces with highly seasonal demands and considers three distinguish strategies as follow: \\
    \begin{center} \scriptsize{
    \begin{tabular}{lrclrclr} \cline{1-2} \cline{4-8}
        & Strategies              &~~~ & Month & Demand &~& Month & Demand   \\ \cline{1-2} \cline{4-8}
      a)& adding seasonal worker  &    &Jan.  & 1600    & & Apr.  & 3800    \\
      b)& using subcontractor     &    &Feb.  & 3000    & & May   & 2200   \\
      c)& building inventory/backlog&  &Mar.  & 3200    & & Jun   & 2200   \\  \cline{1-2} \cline{4-8}
    \end{tabular} }
    \end{center}

    The company sells each tool for \$40 that requires 4 hour and has initial inventory and workers in January of 1000 unit and 80 employees, respectively. A employee of Red Tomato typically works 8 hours a day and earns \$4 per hour during regular time.  Because of labor law, no worker can work more than 20 days of regular time per month and 10 hours of overtime per month.

    \begin{center} \scriptsize{
    \begin{tabular}{lrclrclr} \cline{1-2} \cline{4-5} \cline{7-8}
      Item          & Cost          & ~ & Item        & Cost         & ~ & Item              & Cost\\ \cline{1-2} \cline{4-5} \cline{7-8}
      Material cost & \$15/unit     &   & Hiring costs& \$300/worker &   & RT cost & \$4/hour \\
      Holding cost  & \$2/unit/month&   & Layoff costs& \$500/worker &   & OT cost& \$6/hour \\
      Stockout cost & \$10/unit/month&   & Labor required& 4/unit     &   & Subcontract cost& \$30/unit \\ \cline{1-2} \cline{4-5} \cline{7-8}
    \end{tabular} }
    \end{center}
    Determine the most suitable 'pure' strategy and the optimal 'mixed' strategy that maximizes profit if the company profile must have at least 500 units in June.
  }
  \begin{flushright}
    \tiny{\textbf{adopted from:}  \textsl{Chopra, S. \& Meindl} 2016. pp. 227 }
  \end{flushright}

  \note{
    \textbf{Purpose:} Workshop for aggregate planning for understanding standard policies \\
    \textbf{Leading Question:} How to do aggregate planning? Pure strategies?/Mixed strategies \\
    \begin{itemize}
      \item \textbf{seasonal worker:} fire/hire employees to balance workload (both regular hour and OT)
      \item \textbf{subcontract:} outsource production
      \item \textbf{inventory/backlog:} since demand is fluctuated, why not use inventory and delay shipment to help
    \end{itemize}
  }
\end{frame}


%============================================================================================================
\begin{frame}
  \frametitle{Red Tomato: Model and Excel}

  \begin{center}
    \only<1 |handout:1>{\includegraphics [width=0.85\textwidth]{\figureLocation{scm_demand/redTomatoFlow.png}}}
    \only<2-|handout:0>{\includegraphics [width=0.85\textwidth]{\figureLocation{scm_demand/redTomato.png}}}
  \end{center}

  \note{
    \begin{center}
      \includegraphics [width=0.5\textwidth]{\figureLocation{scm_demand/redTomato.png}}
    \end{center}
    \textbf{Purpose:} Explain MS Excel \\
    \textbf{Leading Question:} How to use MS. Excel to plan?\\
    \begin{itemize}
      \item \textbf{Decision} hire/fire employees, production qty, qty of outsourcing
      \item \textbf{Contraints} shows as OK
      \begin{itemize} \tiny{
        \item \textbf{Using of OT:} must not excess 10 hours per worker
        \item \textbf{Production:} must produce based on OT and regular hour
        \item \textbf{Inventory:} must have the end inventory = 500 }
      \end{itemize}
    \end{itemize}
  }
\end{frame}


%============================================================================================================
\begin{frame}
  \frametitle{Problems with aggregate planning}

  \textsc{Methods}
  \begin{itemize}
    \item \textbf{Price taker:} \alert{price can influence demand} or negotiation (what if) \uncover<2->{$\rightarrow$ S\&OP}
    \item \textbf{Inventory of component:}  lack of components  \uncover<2->{$\rightarrow$ MPS \& MRP}
    \item \textbf{Capacity issues:} simplified production constraints
  \end{itemize}

  \NextLine
  \textsc{Application}
  \begin{itemize}
    \item \textbf{Aggregate product:} no information on exact product \uncover<2->{$\rightarrow$ communication}
    \item \textbf{Delivery plan:} frequency, product mixed, lot size \uncover<2->{$\rightarrow$ communication}
  \end{itemize}

  \note{
    \textbf{Purpose:} limitation of aggregate planning \\
    \textbf{Leading Question:} What we should know about aggregate planning ?\\
    \begin{itemize}
      \item \textbf{Price} Aggre Planning totally ignores qty induced by price
      \item \textbf{Plan $\neq$ Action} production needs components and capacity
      \item \textbf{Capacity} order $\neq$ production now;
      \item[] How to overcome problem in Aggregate Planning
      \item \textbf{Communication:} every concerns can be solved communication
    \end{itemize}

  }
\end{frame}

%============================================================================================================
\section[S\&OP]{Sales and Operations Planning}
%============================================================================================================
\begin{frame}
  \frametitle{HS.0132: Sale and Operation Planning }

  \begin{center}
    \only<1-|handout:1>{\includegraphics [width=0.45\textwidth]{\figureLocation{scm_demand/SOP_concept.png}}}
  \end{center}
   \begin{flushright}
    \tiny{\textbf{source:} ASCM 2020, Principle of S\&OP Workshop }
  \end{flushright}

  \begin{itemize}
    \item \textbf{What:} strategic \alert{supply} or \alert{demand} mgt at \alert{predictable variability}
    \item \textbf{Idea:} seasonal demand, but fixed Cap. \uncover<2->{ $\rightarrow$ aligning finance \& strategic \alert{goals}}
    \item \textbf{Managing supply:} capacity (flex workforce \& machine, subcontracting), investment (common components)
    \item \textbf{Managing demand:} shift or manipulate demand (using price) or promotion
  \end{itemize}

  \note{
    \textbf{Purpose:} Overview on S\&OP\\
    \textbf{Leading Question:} How this price information improve Aggregate Planning  \\
    \begin{itemize}
      \item \textbf{History} P\&G combines demand planning with price and aggregate planning $\rightarrow$ S\&OP
      \item \textbf{managing demands:} non-perishable \& non-fashion,  very seasonal,
      \item \textbf{promotion low demand:} high forward buying, low margin, high holding cost, high switch capacity
      \item \textbf{promotion peak demand:}can steal share, marker expanding, high margin,
      \item[] \textsc{Discussion}
      \item \textbf{chase strategy}
      \item Match the production rate by hiring and laying off employees
      \item Must have a pool of easily trained applicants to draw on
      \item \textbf{Stable workforce} variable work hours
      \item Vary the number of hours worked through flexible work schedules or overtime
      \item \textbf{Level strategy}
      \item Demand changes are absorbed by fluctuating inventory levels, order backlogs, and lost sales
    \end{itemize}
  }
\end{frame}

%================================================================================================-
\begin{frame}
  \frametitle{S\&OP Process }

  \begin{center}
    \only<1-|handout:1>{\includegraphics [width=0.5\textwidth]{\figureLocation{scm_demand/SOP_process}}}
  \end{center}
   \begin{flushright}
    \tiny{\textbf{source:} ASCM 2020, Principle of S\&OP Workshop }
  \end{flushright}

  \begin{itemize}
    \item \textbf{Data Gathering:}  ERP, historical Cap, diff dept EXCEL  \uncover<2->{ $\rightarrow$ \code{current req$^m$}}
    \item \textbf{Demand Planning:} review forecasting, (MTO $\mbox{or}\!\!\!\!?$ MTS) \uncover<2->{ $\rightarrow$ \code{demand plan}}
    \item \textbf{Supply Planning:} last period, involve key suppliers \uncover<3->{ $\rightarrow$ \code{supply plan}}
    \item \textbf{Pre-S\&OP Meeting:} discuss w/ dept on fact, conflict, and plan \uncover<4->{$\rightarrow$ \code{agenda}}
    \item \textbf{Executive Meeting:} review finance and goals for decisions \uncover<5->{$\rightarrow$ \code{S\&OP}}
  \end{itemize}

  \note{
    \textbf{Purpose:} Overview on S\&OP\\
    \textbf{Leading Question:} How this price information improve Aggregate Planning  \\
    \begin{itemize}
      \item \textbf{History} P\&G combines demand planning with price and aggregate planning $\rightarrow$ S\&OP
      \item \textbf{managing demands:} non-perishable \& non-fashion,  very seasonal,
      \item \textbf{promotion low demand:} high forward buying, low margin, high holding cost, high switch capacity
      \item \textbf{promotion peak demand:}can steal share, marker expanding, high margin,
      \item[] \textsc{Discussion}
      \item \textbf{chase strategy}
      \item Match the production rate by hiring and laying off employees
      \item Must have a pool of easily trained applicants to draw on
      \item \textbf{Stable workforce} variable work hours
      \item Vary the number of hours worked through flexible work schedules or overtime
      \item \textbf{Level strategy}
      \item Demand changes are absorbed by fluctuating inventory levels, order backlogs, and lost sales
    \end{itemize}
  }
\end{frame}

%================================================================================================-

\begin{frame}
  \frametitle{Reference}

  \begin{footnotesize}
    \bibliography{\refLocation{Coordination},\refLocation{logisticProvider}}
  \end{footnotesize}

  \note{
    \textbf{Purpose:} overall view of IT in logistics\\
    \textbf{Leading Question:} How each IT system related to one other?\\
    \begin{itemize}
      \item SCM
    \end{itemize}
    [empty] note for citation page
  }

\end{frame}

%=================================================================================================



\end{document}

%=================================================================================================
\section[Exercise]{In-Class Exercise }
%=================================================================================================
\begin{frame}
  \frametitle{FlexMan: Aggre. Planning }

  \footnotesize{
    FlexMan, an electronic contractor, produce tow product categories: routers and switches. Consulting with customers obtains the demands as follows: \\
    \begin{center}
    \begin{tabular}{P{0.07\textwidth}rrr P{0.07\textwidth}rrr P{0.07\textwidth}rr}\hline
      Month & Router & Switch &~& Month & Router & Switch &~& Month & Router & Switch \\ \hline
      Jan.  & 1800   &  1600  & & May   & 800    &  1500  & & Sep.  & 2500   & 1400  \\
      Feb.  & 1600   &  1400  & & Jun   & 1800   &  900   & & Oct.  & 2800   & 1700 \\
      Mar.  & 2600   &  1500  & & Jul   & 1200   &  700   & & Nov.  & 1000   & 800 \\
      Apr.  & 2500   &  2000  & & Aug   & 1400   &  800   & & Dec.  & 1000   & 900  \\ \hline
    \end{tabular}
    \end{center}
    The manufacturing process is govern by numbers of assembly worker, totally 6300 workers. Productions of router and switch require 20 minutes and 10 minutes, respectively. The plant operate 20 days a month, 8 hours a day and paid \$10 per hour and 1.5 time wage for any overtime (maximum 20 hour per worker-month). Total raw materials cost \$20, and handling cost is \$3 per unit-month. At the mean time, the company has no-lay off policy. FlexMan has inventory 100,000 routers and 50000 switches in Jan. The handling costs router and switch are \$2 and \$1 per unit-month.
  }

  \begin{flushright}
    \tiny{\textbf{adopted from:}  \textsl{Chopra  \& Meindl} 2010. pp. 244 }
  \end{flushright}

  \note{
    \begin{itemize}
      \item what is the optimal course of action to fulfil demand
    \end{itemize}
  }
\end{frame}

%============================================================================================================
\begin{frame}
  \frametitle{FlexMan: Questions}

  \small{
    Answer with following questions.
    \begin{itemize}
      \item[1] Assuming no backlog, no subcontractor, and no new hire/leave, what is the optimal production scheduling?
      \item[2] Is any value to increase overtime from 20 hours per worker-month to 40 hour per worker-month
      \item[3] How the solution in (a) and (b) change if the numbers of assembly workers is 5900 and 6700 workers, respectively
      \item[4] Suppose costs of hiring and firing an employee are \$700 and \$1000, respectively. In addition, a new employee has 50\% capacity in her first two months, and then reach full potential later on. What the optimal production scheduling?
      \item[5] If there is a third party charge \$6 per router and \$4 per switch. What the optimal production scheduling?
    \end{itemize}
  }

  \begin{flushright}
    \tiny{\textbf{adopted from:}  \textsl{Chopra, S. \& Meindl} 2010. pp. 244 }
  \end{flushright}

  \note{
    \begin{itemize}
      \item what is the optimal course of action to fulfil demand
    \end{itemize}
  }
\end{frame}


%================================================================================================-
\begin{frame}
  \frametitle{Mass Customization at HP}

  \begin{itemize}
    \item \textbf{Motivation:} Inventory, nature of demand, multiple markets
    \item \textbf{Idea:} \alert{delay the differentiation} of product/process
    \item \textbf{Trade-offs:} economy of scale, production cost \textbf{VS} inventory, transportation
    \item \textbf{Implementation:} cross-organization coordination (manufacturing, purchasing, R\&D, market)
    \item \textbf{Example:} DeskJet, Paint mixture, Benetton shirt, HP hard drive test
    \item \textbf{Achieving mass customization:}
    \begin{itemize}
      \item \code{Ind. product modular:} std. parts, simultaneously manufacturing
      \item \code{Ind. process modular:} flexible machine, separate testing
      \item \code{Agile supply chain:} hired local, reducing permutation
    \end{itemize}
  \end{itemize}

  \note{
    \begin{itemize}
      \item \textbf{Summary:} story of \alert{postponement} HP DeskJet
      \item \textbf{Related concept:} localization, simple and modular design
      \item \textbf{Independent modular product design}
      \begin{itemize}
        \item maximizing standard parts $\rightarrow$ reduce inventory
        \item manufacturing simultaneously $\rightarrow$ reduce leadtime
      \end{itemize}
      \item \textbf{Independent modular process design}
      \begin{itemize}
        \item \textbf{paint mixing machine:} color analyzing machine + base color mixing machine
        \item \textbf{benetton dyeing:} delay dyeing process until realizing customer demands
        \item \textbf{HP hard drive testing process:} standard testing and supplier testing \NextHalfLine
      \end{itemize}
     \item \textbf{Agile supply chain}
      \begin{itemize}
        \item hiring local for value added process (e.g., packaging, assembly, labeling)
        \item reducing inventory and its permutation by postpone assembly process
        \item stock parts and local capacity
      \end{itemize}
    \end{itemize}
  }
\end{frame}

%================================================================================================-
\begin{frame}
  \frametitle{The Games People Play}

  \begin{itemize}
    \item \textbf{Motivation:} tradition forecasting does not work with Games/DVD
    \item \textbf{Reasons:} dense demands in $1^{st}$ period, no history (demands decreasing), \alert{short life cycle}
    \item \textbf{Difficulty:} multiple SKUs (platform, limited edition, package), alternative sources
    \item \textbf{Trade-offs:} dead stock \textbf{VS} out-off-stock
    \item \textbf{Best remedy:} information sharing (no standard protocol), POS, pre-order, search engine info
    \item \textbf{Other examples:} book, popular gadget (e.g., iPhone, iPad, digital camera)
  \end{itemize}

  \note{ \footnotesize{
    \begin{itemize}
      \item Improving accuracy of forecasting by \alert{forecasting parts}, not actual product
      \item Improving leadtime and minimizing disruption by \alert{value added process}
      \item Importance of de-centralizing distribution network,
    \end{itemize}

  \begin{itemize}
    \item \textbf{Trend of products}: fast moving, quick sales, short life cycle, multiple SKUs
    \item[] \hskip3em $\rightarrow$ \alert{dead stock} or \alert{out-off stock}
    \item Improving forecasting by \alert{non-traditional} sources of data
    \item Importance of information exchange and POS data
  \end{itemize}
  }
  }
\end{frame}

%============================================================================================================--
\begin{frame}
  \frametitle{Organizational Planning \forISE}

  \begin{center}
    \only<1 |handout:1>{\includegraphics[width=0.7\linewidth]{\figureLocation{scm_demand/planning/planningCheckList00.png}}}
    \only<2-|handout:0>{\includegraphics[width=0.7\linewidth]{\figureLocation{scm_demand/planning/planningCheckList01.png}}}
  \end{center}
  \begin{flushright}
    \tiny{\textbf{source:} Chase and Jacob. 2011.\hskip10em}\\
  \end{flushright}

  \NextLine
  \textsc{Type of Planning}
  \begin{itemize}
    \item<2-> \textbf{Strategic Planning:} \code{3-5 years} about \code{direction} done by \code{C-Level}
    \item<2-> \textbf{Operation Planning:} \code{1-10 days} about \code{routine} done by \code{staff}
  \end{itemize}


  %%%-------------------
  \note{
    \includegraphics[width=0.45\linewidth]{\figureLocation{scm_demand/planning/planningCheckList01.png}} \\
    \textbf{Purpose:}  \\
    \textbf{Leading Question:} \\
    \begin{itemize}
      \item a
    \end{itemize}
  }
\end{frame}

%============================================================================================================--
\begin{frame}
  \frametitle{General Planning Process \forISE}

  \begin{center}
    \only<1-|handout:1>{\includegraphics[height=0.30\textheight]{\figureLocation{scm_demand/planning/planningProcess.png}}} \\
  \end{center}


  \NextLine
  \begin{block}{How and factor for planing?}{}
    \begin{flushleft} \small{
    \begin{tabular}{lll}
      $\bullet$~any \alert{similar projects} or tasks  & $\bullet$~due date \& consequence & $\bullet$~key resources\\
      $\bullet$~productivity rate (work/FTE)  & $\bullet$~feedback \& review   &   \\
      $\bullet$~risk \& contingency plan      & $\bullet$~visualization     &  \\
    \end{tabular} }
    \end{flushleft}
  \end{block}


  %%%-------------------
  \note{
    \textbf{Purpose:}  \\
    \textbf{Leading Question:} \\
    \begin{itemize}
      \item a
    \end{itemize}
  }
\end{frame}

%============================================================================================================--
\begin{frame}
  \frametitle{Why does planning fail? \forISE}

  \begin{center}
    \begin{quotation}
      He who fails to plan is planning to fail.
    \end{quotation}
  \end{center}
  \begin{flushright}
    \tiny{\textbf{source:} Winston Churchill \hskip8ex ~~~}
  \end{flushright}

  \NextHalfLine

  \begin{tabular}{cl}
    \begin{minipage}{0.30\textwidth}
      \only<1-|handout:1>{\includegraphics[width=\linewidth]{\figureLocation{scm_demand/planning/whyNotPlan.png}}}
    \end{minipage}
  &
    \begin{minipage}{0.65\textwidth}
    \begin{itemize}
      \item Unrealistic \alert{expectations} (too much/ too little)
      \item Lack communication/ engagement/ commitment
      \item Vague responsibility / \alert{misunderstanding}
      \item Unaligned \alert{incentive}/ internal conflict
      \item Unexpected turns (finance/ priority/scope)
      \item Failure to grasp overall planning/ milestone
    \end{itemize}
    \end{minipage}
  \end{tabular}

  %%%-------------------
  \note{
    \textbf{Purpose:}  \\
    \textbf{Leading Question:} \\
    \begin{itemize}
      \item a
    \end{itemize}
  }
\end{frame}
